\documentclass{article}
\usepackage{arxiv}
\usepackage[utf8]{inputenc}
\usepackage[T1]{fontenc}
\usepackage{hyperref}
\usepackage{url}
\usepackage{booktabs}
\usepackage{amsfonts, amsmath}
\usepackage{nicefrac}
\usepackage{microtype}
\usepackage{lipsum}
\usepackage{graphicx}
\usepackage{float}
\usepackage{wasysym} %for special symbols (like gender symbols)

\usepackage{listings}
\usepackage{caption}
\usepackage{xcolor}


%\graphicspath{ {./figures/} }

\title{Structured Data Extraction from Wild-Bee Reports}

\author{
  Felix Böhning
%  \thanks{Use footnote for providing further
%    information about author (webpage, alternative
%    address)}
    \\
  \texttt{felix.boehning@tu-ilmenau.de} \\
  %% examples of more authors
  \And
  Jane Doe \\
  \texttt{jane.doe@tu-ilmenau.de} \\
  %% \AND
  %% Coauthor \\
  %% Affiliation \\
  %% Address \\
  %% \texttt{email} \\
  \AND
  Supervisor: \ldots
}

\begin{document}

\maketitle

\begin{abstract}
\textbf{\textit{Problem:}} The primary challenge addressed in this research is the labor-intensive process of transforming unstructured data into a structured dataset, particularly in the domain of biodiversity research on wild bees. Manual extraction and classification of data from numerous scientific reports is time-consuming and requires deep domain expertise.\linebreak
\textbf{\textit{Questions:}} 

\begin{enumerate}
    \item How effective are open-source Large Language Models (LLMs) like Mistral-7B in automatically extracting structured data from scientific reports on wild bee observations?

    \item How can observations extracted by LLMs be evaluated against manually extracted "gold" labels in a manner that rigorously considers factual accuracy while allowing for variations in expression, given that correct answers can be represented in multiple valid forms?
    
    \item What gains in extraction reliability can be achieved through parameter efficient finetuning (PEFT) on open source LLMs like Mistral-7B?
    
    \item Can synthetically generated samples from a more capable model like GPT-4 improve fine-tuning on smaller open source models?
    
\end{enumerate}

\textbf{\textit{Literature Review:}} The study builds on existing methods for structured data extraction using LLMs. Key techniques include Constrained Prompting, Reasoning, Output-Parsing, and Finetuning, each offering various benefits and limitations in terms of handling diverse data formats and reducing hallucinations. \textbf{\textit{Methodology:}} Mistral-7B-Instruct was finetuned on different dataset combinations including real human annotated reports, synthetic reports with synthetic labels and true negative reports containing no wild bee observations at all. The evaluation pipeline seeks the optimal mapping between labeled observations and predicted observations, identifying false positives and false negatives. It then compares the attributes of true positive observations using a fuzzy score. Precision and recall are subsequently used to derive an overall score. \textbf{\textit{Results:}} Using PEFT, a remarkable performance increase was achieved. The untrained Mistral-7B-Instruct model identified observations with an accuracy of approximately 0.3 and predicted attributes within the true-positive observations with similar accuracy. After three training epochs, observations were extracted with an improved accuracy of approximately 0.8, and attributes also reached an accuracy of approximately 0.8. The use of synthetically generated training data did not yield any noticeable advantages. \textbf{\textit{Conclusions:}} ff

\end{abstract}


\section{Introduction}
Scientific reports predominantly present their information in text form. This poses a challenge when attempting to
compare various reports on a similar topic. Such comparisons are further complicated by data often being presented in
diverse formats, hidden within complex sentence structures, or inferred indirectly from context. Additionally, different
terminologies are used to describe the same concept or phenomenon. Manually extracting and classifying this data
into a consistent tabular structure is labor-intensive. It requires both a deep understanding of the text and expertise in
the relevant field. In biodiversity research on wild bees, the scientific community relies on an extensive collection of
reports, observations, and studies. These records have been contributed over time by both professional apiologists and
amateur biologists and now form the starting point for this work.

Large Language Models (LLMs) showcase impressive capabilities in handling text data, suggesting their use for the
automated extraction of bee observations in structured form. However, LLMs also present the risk of generating
hallucinations, necessitating a critical evaluation of the outcomes from LLM data extraction. Specifically, in specialized
fields like apiology, the general knowledge of LLMs is limited, confining their text comprehension to higher levels of
abstraction.

The objective of this work is to investigate the use of LLMs for the automatic extraction of wild bee observations
in structured form from scientific wild bee reports. An available gold dataset will be used for evaluation of the
extracted data. 


Challenges:
- noisy text
- contextual complexity, draw conclusions from context
- evaluation: natural language --> facts can be expressed in different ways --> allow variability 
    - vs. evaluate factness, detect hallucinations, gabs, mistakes



\section{Related work}

\textbf{Large Language Models (LLMs)} have significantly advanced the field of automatic structured data extraction, introducing several new techniques aimed at improving accuracy and reliability, especially for scientific purposes. However, challenges such as text quality, the reliability of extracted data, and mitigating hallucinations remain central to current research.

\textbf{Text quality and Optical Character Recognition (OCR)} play a crucial role in the success of data extraction processes, particularly when the input text comes from noisy or poorly scanned documents. Methods such as using the Levenshtein distance metric have been employed to evaluate the accuracy of OCR systems and the performance of data extraction \cite{biswas2024robustnessstructureddataextraction}. \cite{10.1145/3639856.3639889} demonstrated that noisy input data can cause rapid degradation of LLM performance in extraction tasks. Additionally, fine-tuning LLMs on noisy datasets often leads to performance issues, as discussed in recent studies \cite{schilling2024text}

\textbf{Prompt engineering} approaches have been instrumental in improving the outputs of large language models (LLMs) for structured data extraction. By providing an initial instruction, the prompts can be iteratively refined to elicit more accurate, relevant, and useful responses tailored to specific use cases. Enriching the prompt with examples facilitates few-shot learning, allowing the model to better understand the expected output. More advanced techniques, such as Chain-of-Thought (CoT) prompting and reasoning, have demonstrated significant improvements in output quality by guiding the model through a structured reasoning process \cite{10.1145/3639631.3639663}. CoT prompting breaks down complex tasks into manageable steps, helping the model reason more effectively through the problem \cite{wei2023chainofthoughtpromptingelicitsreasoning}. Additionally, these approaches allow conversational LLMs to focus on important parts of the input text through conditional prompts and priming, which provides an initial context that directs the model toward a desired output.

\textbf{Reasoning and validation systems} have also been applied to structured data extraction, particularly in fields like medical diagnostics. For example, Bisercic et al. \cite{bisercic2023interpretable} used reasoning combined with value validation systems and handcrafted rules to ensure the extracted data's plausibility. The model employed for medical reports generates structured data in the form of JSON strings, which are then parsed using a pydantic-parser to validate the data’s format and content. Similarly, Li et al. \cite{li2024automatedclinicaldataextraction} introduced a framework for clinical data extraction that aligns internally generated knowledge with external knowledge through in-context learning (ICL). This system employs a retriever to identify relevant knowledge and a grader to evaluate the accuracy and utility of the retrieved data, further improving the reliability of the structured output.

\textbf{Fine-tuning} approaches are another key method used to adapt LLMs to specific tasks, though traditional fine-tuning involves updating all model weights, which is resource-intensive. A more resource-efficient technique is Parameter Efficient Fine-Tuning (PEFT), which has gained popularity for minimizing the number of parameters adjusted during fine-tuning. LoRA (Low-Rank Adaptation) is a widely used PEFT method that adapts less than 1 \% of a model’s parameters while maintaining effectiveness \cite{hu2021loralowrankadaptationlarge}. This approach is particularly valuable because it avoids catastrophic forgetting, thus preserving the model's general abstraction abilities during fine-tuning \cite{zhai2023investigatingcatastrophicforgettingmultimodal}. PEFT has been shown to reduce reliance on human annotations, as demonstrated in the study by Kumar et al. \cite{kumar2024selectivefinetuningllmlabeleddata}, where synthetic labels generated by an LLM were used to reduce manual annotation efforts in schedule-of-event table detection tasks.

\textbf{Human-in-the-loop annotation} is another effective method for balancing the manual labeling effort with the need to ensure reliability, particularly when synthetic labels are prone to hallucinations. This method has been applied to scientific text extraction, where LLMs perform initial annotations, and human experts subsequently correct errors in the annotated data. This corrected data is then used for further fine-tuning, improving the model’s accuracy over time \cite{Dagdelen2024}. 

% LLMs have revolutionized the field of automatic structured data extraction, offering several new techniques. 

% - Recent advances with LLMs enhanced the potential for use in automatic data extraction for scientific purposes

% -main challenges:
%     - text quality, "noisy" text
%     - improve reliability of extracted data and reduction of hallucinations
%     - evaluation

% OCR/TEXT QUALITY:
% - evaluation based on levensthein distance
% - https://arxiv.org/pdf/2406.10295

% - https://dl.acm.org/doi/pdf/10.1145/3639856.3639889
% - finetune with noisy text-data --> performance of llms degrades rapidly

% overview
% - Uni Jena, Material Sciences/ Complete Guide for Data Extraction
% - https://arxiv.org/pdf/2407.16867
% - \ref{schilling2024text} 

% #PROMPT ENGINEERING APPROACHES
% - are trying to guide the LLM to better outputs. 
% - Based on an initial instruction, the prompt can iterativly refined in a way elicits the most accurate, relevant, and useful responses from the model for a specific use-case
% - By enriching the prompt with one or more examples the model can learn (few-shot-learning) more what output is expected
% - More sophisticated prompt engineering like Chain-of-Thought-Prompting and Reasoning lead to further enhancement of the output

% - Chain-of-thought prompting \cite{wei2023chainofthoughtpromptingelicitsreasoning}
% --  decompose the problem into smaller, more manageable steps, and to reason through the problem more effectively

% - A Prompt Engineering Approach for Structured Data Extraction from Unstructured Text Using Conversational LLMs \cite{10.1145/3639631.3639663}
% -- conversational LLMs, chain of thought, reasoning, conditional prompts/ examples
% -- guide the attention of a pre-trained language model (PLM) to the important parts of the input text
% -- Priming provides an initial context to guide the model’s response which can help steer the LLM toward a desired direction or Topic

% - A Simple but Effective Approach to Improve Structured Language Model Output for Information Extraction \cite{li2024simpleeffectiveapproachimprove}
% -- using general knowledge from pretraining
% -- Zero-Shot/ FEw-Shot Learning
% -- Reasoning/ CoT


% - Interpretable Medical Diagnostics with Structured Data Extraction by Large Language Models \cite{bisercic2023interpretable}
% -- difference to Bee Report: medical reports shorter
% --	Reasoning + Value Validation system + handcrafted rules for plausibility (e.g. datatype validation)
% -- extracts structured data as JSON-String that is parsed afterwards using pydantic parser


% - Automated Clinical Data Extraction with Knowledge Conditioned LLMs \cite{li2024automatedclinicaldataextraction}
% --  framework that aligns generated internal knowledge with external knowledge through in-context learning (ICL)
% -- employs a retriever to identify relevant units of internal or external knowledge and a grader to evaluate the truthfulness and helpfulness of the retrieved internal-knowledge rules, to align and update the knowledge bases


% #FINETUNING APPROACHES:
% - in principle adapts all Model-Weights = Continous Pre Training (CPT), but very ressource-consuming
% - A very good approach, which requires minimal ressources is Parameter Efficient Finetuning (PEFT). 

% - LoRA: Low-Rank Adaptation of Large Language Models \cite{hu2021loralowrankadaptationlarge}
% -- is a popular PEFT technique
% -- adapt less than 1 percent of parameters, very effective
% -- no danger of catastrophic forgetting \cite{zhai2023investigatingcatastrophicforgettingmultimodal} -->  preserve the models general abstraction abilities

% - Selective Fine-tuning on LLM-labeled Data May Reduce Reliance on Human Annotation: A Case Study Using Schedule-of-Event Table Detection \cite{kumar2024selectivefinetuningllmlabeleddata}
% -- PEFT Finetuning, used synthetic labels (generted by a LLM) to reduce manual annotation effort


% - Human-in-the-loop annotation = compromise between manual labeling effort vs. guarantee reliability/ avoid hallucinations of synthetic samples

% - Structured information extraction from scientific text with large language models \cite{Dagdelen2024}
% -- Human-in-the-loop annotation: model performs the initial annotation,  human expert corrects possible errors in the annotated data --> further finetuning with corrected data
% -- extracts structured data as JSON-string that is parsed 

\section{Dataset and Features}

% Describe your dataset: how many training/validation/test examples do you have?
% Is there any preprocessing required? 
% What about normalization or data augmentation?
% What is the resolution of your images? 
% How is your time-series data discretized? 
% Include a citation on where you obtained your dataset from. 
% If you plan to extract features using Fourier transforms, word2vec, PCA, ICA, etc. make sure to talk about it. 
% You might also include a representative sample of the dataset. Make sure it is relevant for highlighting your task.

\

- Data given: PDF files and large table with manually extracted attributes


- Quality of PDF files often poor, not all PDFs can be used
- Requirements for PDF Files:
    - symbols must be represented correctly (gender symbols)
    - good OCR Quality
    - no jumbled text
    - content connections must be preserved
- Many PDF files had to be sorted out, because they don't fulfill the requirements

- Variety of attribute formats --> parsing them in one format is not manageable
-- e.g. date: as range, as single day, "summer 1998"
-- e.g. location: coordinates, city, description of forest ways, nature reserve

- attributes representing amounts can not exracted as integers because they are not formatted uniform
-- 



- A uniform data format for extraction is needed
- requirements for the extraction format:
    - unambiguous
    - as compact as possible
    - easy to understand for LLMs/ easy to describe
    - covering as many attributes as possible


- A wild bee observation is clearly described by the scientific name (genus, species, author, year), a location and a time
- We're decided to use the json format for extraction
- To make it more compact similar attributes are concatenated. For example under the key location all location-related attributes (coordinates, TK, city, reserve) are specified
- To make it even more compact redundancies are avoided by using meta attributes. In most cases Observations that are described in the same passage share attributes. Instead of specifing the same attribute for each observation individually it can be specified once as meta-attribute and thereby referes to all subsidary observations.


- For finetuning small handable passages up to 4000 characters should be used
- in a first iteration few short samples where assembled manually by a layman from the few available high quality reports and a model was finetuned on them. After that the model roughly learned the desired output format and was able to perform extractions.

- second step: human in the loop annotation. The pre-finetuned model was integrated into a user interface and wild bee experts were able to extract attributes from their own report passages via online access. The experts were encouraged to critically examine and correct the results, generating high-quality samples for a new finetuning.


- Data-Distribution
    - Number of observations per report passage
    - length of report passage
    - train/test/eval split


\section{ Methods }

% Briefly describe the learning algorithms you plan to use. 
% If there are existing implementations, will you use them and how? 
% How do you plan to improve or modify such implementations?
% Additionally, if you are using a niche or cutting-edge algorithm (anything else not covered in the course), you may want to explain your algorithm using 1/2 paragraphs.


\subsection{Extraction-Pipeline}
- input: report-passage; output: structured json-string representing the described observations with their attributes



- Model should generate a Structured Output as JSON-String using a defined format and attributes
- Extract all attributes as String --> evaluate only strings in evaluation

- each observation is a dictionary with apposite attributes

- save context length:
-- attributes not defined in report passage omitted
-- meta-attribute: if an attribute applies to all observations its only indicated once with prefix 'meta\_' 

see extraction prompt in Appendix \ref{appendix: extraction_prompt}


- Prompting the Model in Conversation-Format (Chat-Format):
-- alternating messages from user, assistant, user, assistant, ... 
-- apply tokenizers chat template which represents conversations as one prompt using special instruction tokens (depending on model)

< A figure that displays the conversation >

- simulating conversation context. Initial prompt by user, assistant replies with 'okay lets start' (necessary to keep the alternating-structure), three few shot examples with gold completion... only last prediction is made by model

- Finetuning
-- apply MCE-Loss only between last prediction and gold completion 


\subsection{Evaluation-Pipeline}


% Define custom style for JSON formatting with proper coloring
\lstdefinelanguage{json}{
    basicstyle=\ttfamily\scriptsize, % Smaller text size
    stringstyle=\color{brown},
    keywordstyle=\color{blue},
    commentstyle=\color{gray},
    numbers=none,
    showstringspaces=false,
    breaklines=true,
    morestring=[b]",
    morekeywords={true, false, null},
    columns=flexible, % Makes the text narrower
    literate=%
      {ä}{{\"a}}1 {ö}{{\"o}}1 {ü}{{\"u}}1 % Handle ä, ö, ü
      {Ä}{{\"A}}1 {Ö}{{\"O}}1 {Ü}{{\"U}}1 % Handle Ä, Ö, Ü
      {ß}{{\ss}}1     
}


\begin{figure}[h]  % Add [h] to place it "here"
  \noindent

  \begin{minipage}{0.49\textwidth}
      \centering
      Prediction \\ % Header above JSON 1
      %\vspace{0.3cm} % Space between header and JSON
      \lstset{language=json}
      \begin{lstlisting}
      {
        "meta_scientific_name": "Andrena fucata (SMITH 1847)",
        "meta_location": "Nordrhein-Westfalen, Bielefeld",
        "meta_collecting_method": "Streifnetz (Handfang) & Totfunde",
        "observations": [
          {
            "date": "1993, 1996",
            "location": "Universität"
          },
          {
            "date": "1996",
            "location": "Dornberger Str."
          }
        ]
      }
      \end{lstlisting}
  \end{minipage}
  \hfill
  \begin{minipage}{0.49\textwidth}
      \centering
      Label (Gold) \\ % Header above JSON 2
      %\vspace{0.0cm} % Space between header and JSON
      \lstset{language=json}
      \begin{lstlisting}
      {
        "meta_scientific_name": "Andrena fucata (SMITH 1847)",
        "meta_date": "1993, 1996",
        "meta_location": "Nordrhein-Westfalen, Bielefeld",
        "meta_collecting_method": "Streifnetz (Handfang) & Totfunde",
        "observations": [
          {
            "location": "Universität"
          },
          {
            "location": "Dornberger Str."
          }
        ]
      }
      \end{lstlisting}

  \end{minipage}
  \noindent
  \captionsetup{justification=raggedright,singlelinecheck=false} % Left-align caption
  \captionof{figure}{Example for predicted- and labeled Observation-Extraction}
  \vspace{0.5cm} % Space between the JSONs and caption
\end{figure}


\lstdefinestyle{pseudocode}{
    basicstyle=\ttfamily\small, % Use monospaced font for code
    columns=fullflexible,
    mathescape=true, % Allows escaping to LaTeX math mode within $
    keepspaces=true,
    numbers=left, % Line numbers on the left
    numberstyle=\tiny\color{gray}, % Line numbers styling
    stepnumber=1, % Line number step
    numbersep=5pt, % How far the line numbers are from the code
    backgroundcolor=\color{white}, % Background color for the code block
    showspaces=false, % Don't show space characters
    showstringspaces=false, % Don't show spaces in strings
    showtabs=false, % Don't show tabs
    frame=single, % Adds a frame around the code
    tabsize=2, % Sets default tabsize to 2 spaces
    captionpos=b, % Sets the caption-position to bottom
    breaklines=true, % Set automatic line breaking
    breakatwhitespace=false, % Sets if automatic breaks should only happen at whitespace
    escapeinside={(*@}{@*)}, % if you want to add LaTeX within your code
    morecomment=[l]{//}, % Comment style
    commentstyle=\color{gray}\ttfamily,
    xleftmargin=\parindent,
    keywordstyle=\bfseries % Bold keywords
}


\begin{figure}[ht]
    \centering
    \begin{lstlisting}[style=pseudocode]
    # Step 1: Parse predicted and gold JSON strings
    pred_json = parse_json(pred_json_string)
    gold_json = parse_json(gold_json_string)
    
    # Step 2: Compare observations with matching scientific names
    for pred_obs in pred_json:
        for gold_obs in gold_json:
            if pred_obs['scientific_name'] == gold_obs['scientific_name']:
                attribute_scores = []
    
                # Step 3: Compare attributes between the matched observations
                for attr in all_attributes:
                    if attr in pred_obs and attr in gold_obs:
                        score = fuzzy_ratio(pred_obs[attr], gold_obs[attr])
                    else:
                        score = 0  # Attribute is missing or hallucinated
                    attribute_scores.append(score)
    
                # Step 4: Calculate mean similarity for this pair
                sim_score = mean(attribute_scores)
                observation_pairs.append((pred_obs, gold_obs, sim_score))
    
    # Step 5: Sort pairs by similarity score
    observation_pairs.sort(key=lambda x: x[2], reverse=True)
    
    # Step 6: Map the highest similarity pairs (ensure uniqueness)
    used_pred, used_gold = set(), set()
    tp_pairs = []
    for pred_obs, gold_obs, _ in observation_pairs:
        if pred_obs not in used_pred and gold_obs not in used_gold:
            tp_pairs.append((pred_obs, gold_obs))
            used_pred.add(pred_obs)
            used_gold.add(gold_obs)
    \end{lstlisting}
    \caption{Pseudocode for Observation-Mapping Algorithm}
    \label{fig:mapping_pseudocode}
\end{figure}

% - Challenge: Compare predicted json-string with gold-json-string

% <a figure displaying a predicted observation json and a the respective label>

% - Why not use MCE-Loss for evaluation?
% -- MCE compares prediction and gold tokens per position, but corresponding attributes/ values can be located in different positions
% -- MCE can only be calculated between existing predictions and gold-tokens, so hallucinated attributes (False Positives) or unrecognized attributes (False Negatives) can't be included in evaluation

% - Required: Mapping between predicted observations and gold observations
% -- first-step: parsing both json strings (prediction and gold) to json, normalize observations by assigning the meta-attributes

% -- for all observations concerning bees with matching scientific-name:
% --- for all attributes:
% ---- if attribute specified in predicted observation and gold observation --> fuzzy-ratio (based on Levenshtein-Distance \cite{fuzzywuzzy2023}) --> Score between 0 (worst)... 1 (best, exact match)
% ---- if attribute specified in predicted observation and not in gold observation --> assume hallucination (FP) --> score 0
% ---- if attribute not specified in predicted observation but in gold observation --> assume not recognized (FN) --> score 0
% --- calculate the mean-score over all attributes for this observation pair = attribute-similarity


% -- now every pair of predicted and gold observation concerning the same bee (scientific name) has a score based on their attribute-similarity 
% -- order the pairs based on their attribute-similarities
% -- collect True Positive (TP) Observation-Pairs by going through the ordered list beginning with the Pair that has the highest attribute-similarity 
% --- if an Observation-Pair consists of a predicted or gold-observation that was already used 
% --- achieve uniqueness: if a predicted observation or a gold-observation was already involved in a TP-Pair it is ignored

% -- thus finding the best similarity-mapping between predictions and gold-observations

% -- identify hallucinated observations (False Positives): predicted observation that are not involved in any TP-Pair
% -- identify not-recognized observations (False Negatives): gold observations that are not involved in any TP-Pair

% <mapping algorithm figure>

% Based on the number of TP-, FP- and FN- Observations we now can calculate Precision (P), Recall (R) and F1-Score (F1) describing how well the model was able to identify observations in the report-passage.
% - thereby attributes only involved implicitly
% -- problem: its a greedy mapping.. if for one bee (scientific name) same number of predicted and gold observations they were always completely mapped as TP-Pairs, even if their attribute-similarity is very low
% -- Solution: include attribute-scores

% - inclusion of attributes into evaluation
% -- introducing two different kinds of True Positives:
% --- TP: number of attributes specified in both observations, regardless of their values 
% --- TP\textsubscript{fuzzy}: the sum of fuzzy-ratios between all attributes specified in both observations, if attributes match exactly: TP\textsubscript{fuzzy} equals TP

% \begin{equation}
% P\textsubscript{fuzzy} = \frac{TP\textsubscript{fuzzy}}{TP+FP};\quad \quad
% R\textsubscript{fuzzy} = \frac{TP\textsubscript{fuzzy}}{TP+FN};\quad \quad  
% F1\textsubscript{fuzzy} = \frac{2*P\textsubscript{fuzzy}*R\textsubscript{fuzzy}}{P\textsubscript{fuzzy}+R\textsubscript{fuzzy}}
% \label{eq:attribute_score}
% \end{equation}

% - By combining F1 of Observations with F1\textsubscript{fuzzy} of attributes receive a total score
% \begin{equation}
% S = F1\textsubscript{observations}*F1\textsubscript{fuzzy,attributes} \textrm{ [0: worst ... 1: best]}
% \label{eq:total_Score}
% \end{equation}

\textbf{Mapping Predictions to Gold Observations:} The primary challenge lies in comparing predicted observations in the form of JSON strings with their gold-standard counterparts. A direct comparison using standard metrics such as Mean Cross Entropy (MCE) proves inadequate due to the inherent variability in the positioning of attributes within the JSON structure. While MCE compares token sequences at fixed positions, it cannot account for discrepancies such as hallucinated attributes (false positives) or missing attributes (false negatives).

To address these issues, we introduce a mapping procedure that pairs predicted and gold observations based on matching scientific names of bee species. The JSON strings representing these observations are first parsed and normalized, ensuring consistent meta-attribute assignment across both predicted and gold-standard data.

For each pair of observations with matching scientific names, we compare their attributes as follows:

\begin{itemize}
    \item If an attribute is present in both the predicted and gold observation, we compute a fuzzy similarity score using the Levenshtein distance \cite{fuzzywuzzy2023}, ranging from 0 (no match) to 1 (exact match).
    \item If an attribute exists in the predicted observation but not in the gold-standard, we consider it a hallucination (false positive) and assign a score of 0.
    \item Conversely, if an attribute is present in the gold observation but absent in the prediction, we consider it not recognized (false negative), also scoring 0.
\end{itemize}

The mean similarity score across all attributes in an observation pair is referred to as the attribute similarity score. \\

\textbf{Identifying True Positives, False Positives, and False Negatives:} Each predicted observation is paired with the gold observation with the highest attribute similarity score, creating a list of potential true positive pairs. We enforce uniqueness in this pairing process, ensuring that each observation is used only once. Unpaired predicted observations are classified as hallucinations (false positives), and unpaired gold observations as unrecognized (false negatives).

This greedy mapping strategy allows us to establish a best-match similarity between predicted and gold-standard observations. However, a limitation of this approach is that it may result in true positive pairs with low attribute similarity when the number of observations matches exactly. To mitigate this, we propose an additional step that incorporates attribute scores into the evaluation.\\


\textbf{Attribute-Level Evaluation:} To refine our evaluation further, we introduce two types of true positives:

\begin{itemize}
    \item TP: The count of attributes that are present in both the predicted and gold observations, irrespective of their values.
    \item TP\textsubscript{fuzzy}: The sum of fuzzy similarity scores between all matching attributes. If attributes are an exact match, TP\textsubscript{fuzzy} equals TP.
\end{itemize}


Using these metrics, we define precision, recall, and F1-score at the attribute level as follows:

\begin{equation}
    P\textsubscript{fuzzy} = \frac{TP\textsubscript{fuzzy}}{TP+FP};\quad \quad 
    R\textsubscript{fuzzy} = \frac{TP\textsubscript{fuzzy}}{TP+FN};\quad \quad
    F1\textsubscript{fuzzy} = \frac{2*P\textsubscript{fuzzy}*R\textsubscript{fuzzy}}{P\textsubscript{fuzzy}+R\textsubscript{fuzzy}}
    \label{eq:pre_rec_f1_fuzzy} 
\end{equation}\\

\textbf{Final Evaluation Score:} We combine the F1-score of observations with the F1\textsubscript{fuzzy} of attributes to calculate a total evaluation score that accounts for both observation-level and attribute-level performance:

\begin{equation}
    S = F1\textsubscript{observations}*F1\textsubscript{fuzzy,attributes} \textrm{ [0: worst ... 1: best]} 
    \label{eq: total_score} 
\end{equation}

This approach provides a comprehensive evaluation of the model’s ability to extract structured data, reflecting its performance at both the observation and attribute levels.


\section{Experiments/Results/Discussion}
% You should also give details about what (hyper)parameters you chose, e.g., why did you use X learning rate for gradient descent, what was your mini-batch size and why, and how you chose them. 
% What your primary metrics are: accuracy, precision, AUC, etc. 
% For results, you want to have a mixture of tables and plots.
% If you are solving a classification problem, you should include a confusion matrix or AUC/AUPRC curves. 
% Include performance metrics such as precision, recall, and accuracy. 
% For regression problems, state the average error. 
% You should have both quantitative and qualitative results. 
% If it applies: include visualizations of results, heatmaps, examples of where your algorithm failed and a discussion of why certain algorithms failed or succeeded. 
% In addition, explain whether you think you have overfit to your training set and what, if anything, you did to mitigate that. 
% Make sure to discuss the figures/tables in your main text throughout this section. 
% Your plots should include legends, axis labels, and have font sizes that are legible when printed


\section{Experiments}

in all experiments:
- Model used: meta-llama/Meta-Llama-3.1-8B-Instruct
<overview of model loading parameters>
- all model interaction carried out in chat-mode applying the tokenizers chat template
    - user provides prompt and report passage
    - assistant answers with json string


Inference setup:
- break up generation after json is completed (\} ```)
- <overview of generation params>

 
Finetuning setup:
- same lora configuration in all experiments 
- number/ percentage of parameters trained in this lora setting!
<overview of lora parameters>

- finetuning using mce-loos for backpropagation
- calculate loss only on completion (last assitant message)


Three Major Finetuning-Experiments, each carried out over 10 training-epochs:
\begin{enumerate}
    \item \textbf{prompt combined with 3 few-shot examples:} The prompt is the first user message. Because user and assistant messages must be alternating in this case the assistant answers with a short 'ok, lets go' followed by the first user-provided few-shot-example passage, followed by the belonging assistant answer. Only the answer to the last user provided report passage is really generated and included into the loss calculation
    \item \textbf{only prompt:} 
    \item \textbf{without prompt:}
\end{enumerate}

- In each experiment the unfinetuned base model will be evaluated as baseline. After that evaluation is conducted after each epoch.

- Using a prompt combined with few-shot samples inducts more context than just providing a single report-passage without explanation or samples. In terms of scalability self-attention scales quadratically with the context length in terms of both computation and memory. On the other hand it enforces 'in-context-learning' where the model faster understands the task from standing. Therefore we expect a faster converging model in the first two finetuning-experiments

- Regarding the third experiment we expect the model to converge slower. Its going to be interesting to see if the model can adapt to the extraction task without further explanation. The model is required to learn which attributes are used just indirectly from the loss backpropagation labels. If the results are good, this offers a good opportunity to save context-length and thereby inputting longer report-passages in production use.

- It should also considered to use 'best of both worlds', by finetuning the model with a context-rich fast converging method and later do the inference with less additional context as possible.


\section{Results and Discussion}

< a figure comparing the experimental setups>




\section{Conclusion/Future Work}


\section{Contributions}
This section should describe what each team member worked on and contributed to the project.
The contributions section is not included in the 5 page limit. 

\appendix
\section{Extraction-Prompt} \label{appendix: extraction_prompt}


\begin{figure}[H]
    \centering
    \fbox{%
        \parbox{0.95\textwidth}{
            %\raggedright % Makes the content left-aligned
            %\raggedright % Makes the content left-aligned
            \ttfamily % Monospaced font
            %\fontfamily{cmtt}\selectfont % Use a monospaced font
            \fontsize{8pt}{10pt}\selectfont % Set font size
            %\setlength{\baselineskip}{8pt} % Adjust line spacing here
            Your task is to extract data on wild bee observations described in the report passages into structured JSONs. The expected json format is specified below. Extract only the specified attributes. If an attribute of an observation is not specified or you can't find information for this attribute, leave it out. Don't make up information! Stick strictly to the information in the report passage. Extract information exactly as it appears in the passage.\\

            ATTRIBUTE DESCRIPTIONS:\\
            bee: Latin designation of the bee. Consisting of genus and species\\
            taxon: name of the author who first described the bee, and sometimes the respective publication year. In most cases the taxon is located right behind the latin designation of the bee\\
            n\_males: number of observed male bees, could be abbreviated with "\male"-symbol after the respective amount\\
            n\_females: number of female bees/ worker bees, could be abbreviated with "\female"-symbol after the respective amount\\
            n\_queens: the number of queens is explicitly described\\
            n\_divers: number of bees with unassignable gender, could be abbreviated with "\mercury"-symbol after the respective amount\\
            leg: person who confirmed the bee, could be abbreviated with "leg." in front of the respective name\\
            coll: person who collected the bee, could be abbreviated with "coll." in front of the respective name\\
            det: person who has determined the bee, could be abbreviated with "det." or "vid." in front of the respective name\\

            All attributes in their defined order:\\
            bee, taxon, date, location, n\_males, n\_females, n\_queens, n\_divers, leg, coll, det, habitat, visited\_flowers, behaviour, observed\_nesting, collecting\_method, remarks\\
            
            EXPECTED JSON FORMAT:\\
            
            The JSON shall contain a list of observations as value for the key "observations". Each observation should be a dictionary containing the sorted attributes specified above as keys and observation-specific values extracted from the report passage. If an attribute is not specified for an observation, leave it out. If an attribute has the same value for all observations, it should be reported as meta-attribute with the prefix "meta\_". Meta-attributes should be specified before "observations" in the order defined above. If an attribute is specified as a meta-attribute, it should not be repeated in the observation-dictionary. Only "location" can be specified as meta-attribute and observation-specific attribute at the same time. If you can't find any wild bee observations in the report passage return an empty JSON.}
        }%
    \caption{Extraction-Prompt consisting of initial task instruction, explanation of attributes, order of further attributes and description of the output-format}
\end{figure}


%\section*{References}

\bibliographystyle{unsrt}  
\bibliography{references} 

\end{document}


